\documentclass{article}
\usepackage[utf8]{inputenc}
\usepackage[T1]{fontenc}
\usepackage{amsmath}
\usepackage{amsfonts}
\usepackage{booktabs}
\usepackage{listings}
\usepackage{xcolor}
\usepackage{geometry}

\geometry{margin=1in}

\title{Programming Coursework Summary}
\author{Disa Degerholm}
\date{\today}

% Color definitions for Python syntax highlighting
\definecolor{codegreen}{rgb}{0,0.6,0}
\definecolor{codegray}{rgb}{0.5,0.5,0.5}
\definecolor{codepurple}{rgb}{0.58,0,0.82}
\definecolor{backcolour}{rgb}{0.95,0.95,0.92}

\lstset{
    language=Python,
    backgroundcolor=\color{backcolour},
    commentstyle=\color{codegreen},
    keywordstyle=\color{magenta},
    numberstyle=\tiny\color{codegray},
    stringstyle=\color{codepurple},
    basicstyle=\ttfamily\small,
    breakatwhitespace=false,
    breaklines=true,
    captionpos=b,
    keepspaces=true,
    numbers=left,
    numbersep=5pt,
    showspaces=false,
    showstringspaces=false,
    showtabs=false,
    tabsize=4
}

\begin{document}

\maketitle

\section{Code Style \& Layout}

\subsection{Indentation \& Spacing}
\begin{itemize}
    \item \textbf{Indentation:} Always use exactly 4 spaces per indentation level without mixing tabs. 
    \item[] \textit{Example:} If statements must indent their nested logic block by exactly four empty spaces.
    \item \textbf{Operators:} Place a single space on both sides of assignment, mathematical, and comparison operators.
    \item[] \textit{Example:} Write \texttt{x = y + 2} instead of \texttt{x=y+2}.
    \item \textbf{Brackets:} Do not leave spaces immediately inside parentheses or brackets, nor between function names and brackets.
    \item[] \textit{Example:} Write \texttt{print("test")} and \texttt{items[0]} tightly without internal padding spaces.
\end{itemize}

\subsection{Blank Lines}
\begin{itemize}
    \item \textbf{Two Blank Lines:} Use exactly two empty rows to separate top-level functions and classes.
    \item[] \textit{Example:} Place two blank lines down before declaring a new standalone \texttt{def global\_func():}.
    \item \textbf{One Blank Line:} Use a single empty row to separate methods inside a class or distinct logic blocks inside a function body.
    \item[] \textit{Example:} Insert one blank line between the definition loops inside a calculation function block.
\end{itemize}

\section{Definitions \& Syntax Logic}

\subsection{Environment Setup}
Intuition: Pass the active text configuration to the interpreter by pressing F5 or the green triangle.
\begin{lstlisting}[language=Python]
print("Welcome to DA2005")
\end{lstlisting}

\subsection{Variables \& Assignment Logic}
Intuition: Variables act as named boxes where evaluation flows right-to-left.
\begin{lstlisting}[language=Python]
distance = 67
time_used = 1.25
average_velocity = distance / time_used
\end{lstlisting}

\subsection{Print Separators}
Intuition: Commas inside execution brackets natively inject a space separation.
\begin{lstlisting}[language=Python]
print("Velocity:", average_velocity)
\end{lstlisting}

\subsection{Multi-Variable Assignment}
Intuition: Commas separate parallel variable groups to match sequence layouts linearly.
\begin{lstlisting}[language=Python]
x, y = 2, 3
\end{lstlisting}

\subsection{Division Systems}
Intuition: Single slash outputs a decimal float, double slash truncates to an integer.
\begin{lstlisting}[language=Python]
z = 22 / 7
q = 22 // 7
\end{lstlisting}

\subsection{Exponentiation Operator}
Intuition: Double up the asterisk multiplication symbol to signify power scaling.
\begin{lstlisting}[language=Python]
power = 2 ** 3
\end{lstlisting}

\subsection{Modulo Operation}
Intuition: The percentage symbol returns only the leftover remainder from a division.
\begin{lstlisting}[language=Python]
remainder = 2 % 3
\end{lstlisting}

\subsection{Alternating Quotes}
Intuition: Swap double and single enclosing symbols to print text quotes cleanly.
\begin{lstlisting}[language=Python]
quote_1 = "Sa du 'hej'?"
\end{lstlisting}

\subsection{String Escaping}
Intuition: Embedding a backslash before an n pushes printing down to a newline.
\begin{lstlisting}[language=Python]
multiline_print = "Hipp\nnhipp\nnhurra!"
\end{lstlisting}

\subsection{User Interactivity}
Intuition: Halts operations, displays the prompt text, and returns a string object.
\begin{lstlisting}[language=Python]
namn = input("Vad heter du?\n")
\end{lstlisting}

\subsection{Robust String Searching}
Intuition: Map a letter to a number safely using find; missing values yield -1.
\begin{lstlisting}[language=Python]
def hexa_char_to_int(symbol: str) -> int:
    hexa_chars = "0123456789ABCDEF"
    return hexa_chars.find(symbol.upper())
\end{lstlisting}

\subsection{Native Brackets Lookup}
Intuition: Map a number to a letter by passing integer indices into square brackets.
\begin{lstlisting}[language=Python]
def int_to_hexa_char(tal: int) -> str:
    hexa_chars = "0123456789ABCDEF"
    return hexa_chars[tal]
\end{lstlisting}

\subsection{Guard Clauses}
Intuition: Escape validation loops early with condition gates to prevent crashes.
\begin{lstlisting}[language=Python]
def deci_to_hex(deci: int) -> str:
    if deci == 0:
        return "0"
\end{lstlisting}

\section{Good Practice Ideas}

\begin{itemize}
    \item \textbf{English Preference:} Standardize naming, comments, and project documentation in English.
    \item[] \textit{Example:} Name an age variable \texttt{user\_age} instead of using Swedish terms.
    \item \textbf{Self-Documenting Code:} Assign meaningful identities to identifiers so logic states remain obvious.
    \item[] \textit{Example:} Use \texttt{average\_velocity} rather than the ambiguous character \texttt{v}.
    \item \textbf{Inline Comments:} Leverage hash marks to declare why specific strategies are applied.
    \item[] \textit{Example:} Write \texttt{\# Safely handles mathematical limits} above boundary conditions.
    \item \textbf{Literal Management:} Isolate hardcoded constants at the script top for fast adjustments.
    \item[] \textit{Example:} Declare \texttt{TAX\_RATE = 0.25} globally instead of inserting 0.25 within expressions.
    \item \textbf{Type Constraints:} Explicitly convert input outcomes since they are captured as text strings.
    \item[] \textit{Example:} Cast numerical input using \texttt{age = int(input("Age: "))}.
    \item \textbf{Reserved Word Auditing:} Avoid reassigning internal keywords or core tool tags.
    \item[] \textit{Example:} Never declare a tracking variable with native system names like \texttt{if = 4}.
    \item \textbf{Guard Clauses:} Deal with anomalous inputs instantly using standalone control blocks.
    \item[] \textit{Example:} Return zero instantly if checking collections containing zero entities.
    \item \textbf{Boolean Tests:} Evaluate state queries directly through native logical conditions.
    \item[] \textit{Example:} Validate flags with syntax structured as \texttt{if not ready:}.
    \item \textbf{String Searching:} Swap string index calls with find methods to suppress runtime value errors.
    \item[] \textit{Example:} Execute \texttt{"ABC".find("Z")} to gather a manageable negative integer flags.
    \item \textbf{Typo Prevention:} Maintain method verification practices to ensure file streams release memory.
    \item[] \textit{Example:} Ensure text handles clear active locks explicitly via the \texttt{.close()} routine.
    \item \textbf{Interactive Safeguards:} Loop defensive error blocks to secure input pipelines.
    \item[] \textit{Example:} Run user queries inside continuous checks protected by try and except blocks.
\end{itemize}

\section{Project Documentation \& Git}

\subsection{Markdown Layout}
\begin{itemize}
    \item \textbf{Structure:} Define explicit primary headers above descriptive overview paragraphs.
    \item[] \textit{Example:} Write \texttt{\# Project Name} at the absolute first line of text layout files.
    \item \textbf{Compliance Checklists:} Record architectural limitations using bullet points to demonstrate obedience.
    \item[] \textit{Example:} Include a checklist item stating that built-in libraries were avoided.
    \item \textbf{Code Presentation:} Bind technical scripts inside specialized language blocks using triple backticks.
    \item[] \textit{Example:} Wrap programming examples within block regions declared as \texttt{\textasciigrave\textasciigrave\textasciigrave python}.
\end{itemize}

\subsection{Git Version Control}
\begin{itemize}
    \item \textbf{Repository Functionality:} Think of repositories as secure project historical records tracking file adjustments.
    \item[] \textit{Example:} Use repository history logs to roll back accidental formatting changes.
    \item \textbf{Public Visibility:} Deselect visibility restrictions during cloud uploads to provide structural links.
    \item[] \textit{Example:} Toggle cloud visibility off from private mode inside your local desktop sync application.
\end{itemize}

\subsection{Comparison Matrix}
\begin{table}[h]
\centering
\begin{tabular}{@{}lll@{}}
\toprule
\textbf{Feature} & \textbf{Normal Desktop Folder} & \textbf{GitHub Repository} \\ \midrule
Storage & Lives locally on hard drives. & Stored locally and backed up online. \\
Version History & Overwritten edits are lost. & Tracks every historical snapshot taken. \\
Sharing & Requires file compression and emailing. & Shared instantly via a public web URL. \\ \bottomrule
\end{tabular}
\caption{Desktop Folder vs. GitHub Repository}
\end{table}

\end{document}